\subsection{Maximal Étale Subalgebras}

Lemma 4. --- Let A be a central simple algebra of finite degree over K, distinct from K. There exists an étale (V, §6, No.~3, p.~28, Definition 1) subalgebra of A distinct from K.

By Wedderburn's theorem (VIII, p.~120, Theorem 1), we may assume that \(\mathrm{A}\) is of the form \(\mathbf{M}_n(\mathrm{D})\) , where \(n\) is a strictly positive integer and \(\mathrm{D}\) a field with center \(\mathrm{K}\) .

Suppose n \textgreater{} 1. The algebra of diagonal matrices with entries in K is an étale subalgebra of A distinct from K.

Suppose n = 1. Let p be the characteristic exponent of D. By Lemma 1 of VIII, p.~230, there exists an element a of D such that \(a^{p^{m}}\) does not belong to K for any natural number m. So for m sufficiently large, the element \(x = a^{p^{m}}\) is separable over K (V, §7, No.~7, p.~44, Proposition 13), but it does not belong to K. The subalgebra \(\mathrm{K}(x)\) of A is a separable extension of the field K of finite degree, hence an étale subalgebra over K; it is distinct from K.

PROPOSITION 4. --- Let A be a central simple algebra of finite degree over K. Let L be a subalgebra of A, and let L' be the commutant of L in A.

\begin{enumerate}
\def\labelenumi{\alph{enumi})}
\item
  If L is maximal among the semisimple commutative subalgebras of A, then we have \(L = L'\) , and L is a maximal commutative subalgebra of A.
\item
  If L is maximal among the étale subalgebras of A, then we have \(L = L'\) , and L is a maximal commutative subalgebra of A.
\end{enumerate}

We know that the relation \(L = L'\) means that L is a maximal commutative subalgebra of A (VIII, p.~261, Lemma 3, a)). Suppose that L is semisimple, commutative, and distinct from \(L'\) . By Theorem 5 of VIII, p.~259, \(L'\) is semisimple, and L is the commutant of \(L'\) and therefore the center of \(L'\) ; consequently, \(L'\) is not commutative. It suffices to prove that there exists a semisimple commutative subalgebra M of A that is distinct from L and contains L and that is étale if L is étale.

By the structure theorem for semisimple rings (VIII, p.~135, Theorem 1), there exist simple rings \(B_{1},\ldots,B_{r}\) and an isomorphism \(\varphi\) from \(L'\) to \(B_{1}\times\cdots\times B_{r}\) . For \(1\leqslant i\leqslant r\) , denote the center of \(B_{i}\) by \(E_{i}\) ; we then have \(\varphi(\mathrm{L})=\mathrm{E}_{1}\times\cdots\times\mathrm{E}_{r}\) . Since \(L'\) is not commutative, we may assume that \(B_{1}\) , for example, is not commutative; we therefore have \(B_{1}\neq E_{1}\) , and, by Lemma 4, there exists a subalgebra \(M_{1}\) of \(B_{1}\) that is commutative, distinct from \(E_{1}\) , and étale over \(E_{1}\) . Set \(\mathrm{M} = \varphi^{-1}(\mathrm{M}_{1} \times \mathrm{E}_{2} \times \cdots \times \mathrm{E}_{r})\) ; it is a semisimple commutative subalgebra of A containing L and distinct from L. Suppose that L is étale over K; let us prove that M is étale. The extensions \(E_{i}\) of K are separable (V, §6, No.~4, p.~30, Proposition 3). Moreover, since the \(E_{1}\) -algebra \(M_{1}\) and the K-algebra \(E_{1}\) are étale, the K-algebra \(M_{1}\) is also étale (V, §6, No.~5, p.~32, Corollary 2). Thus the K-algebra \(M_{1} \times E_{2} \times \cdots \times E_{r}\) is étale, and therefore so is M.

Let A be a central simple algebra of finite degree over K. A subalgebra of A that is maximal among the semisimple commutative subalgebras of A is called a maximal semisimple commutative subalgebra of A. By Proposition 4, the term ``maximal'' refers to the property of being commutative or of being semisimple and commutative. A subalgebra of A that is maximal among the étale subalgebras of A is called a maximal étale subalgebra of A.

COROLLARY 1. --- Let A be a central simple K-algebra of finite degree. Every semisimple (resp. étale) commutative subalgebra of A is contained in a maximal commutative subalgebra of A that is semisimple (resp. étale).

COROLLARY 2. --- Let D be a field of finite degree over K with center K.

\begin{enumerate}
\def\labelenumi{\alph{enumi})}
\item
  The maximal commutative subfields of D are the maximal commutative subalgebras of D and also the maximal semisimple commutative subalgebras of D. Every commutative subfield L of D is contained in a maximal commutative subfield.
\item
  Let L be a commutative subfield of D that is a separable extension of K; it is contained in a maximal commutative subfield of D that is a separable extension of K.
\item
  Let L be a commutative subfield of D. Then L is a maximal commutative subfield of D if and only if we have \([D:K]=[L:K]^{2}\) .
\end{enumerate}

A subalgebra of D is a field (V, §2, No.~2, p.~10, Proposition 1) and is therefore semisimple. Moreover, a maximal commutative subfield of D contains K. Assertions a) and b) then follow from Corollary 1 and assertion c) of Proposition 3 (VIII, p.~262).

PROPOSITION 5. --- Let A be a central simple algebra of finite degree over K. Let B be a semisimple subalgebra of A, and let B' be the commutant of B.

\begin{enumerate}
\def\labelenumi{\alph{enumi})}
\item
  The algebra B contains a maximal semisimple commutative subalgebra of A if and only if B contains B'.
\item
  Suppose that B contains \(B'\) , and let g be a K-algebra homomorphism from B to A. There exists an inner automorphism \(\theta\) of A that coincides with g on B.
\end{enumerate}

Let L be a maximal commutative subalgebra of A; by Lemma 3 of VIII, p.~261, L is equal to its commutant \(L'\) in A. If B contains L, then its commutant \(B'\) is contained in \(L'\) and therefore in B.

Conversely, suppose that \(B'\) is contained in B. Then \(B'\) is the center of B, and it is a semisimple commutative algebra (VIII, p.~137, Proposition 2). By Corollary 1, there exists a maximal semisimple commutative subalgebra L of A containing \(B'\) . The commutant of L is L (VIII, p.~261, Lemma 3, a)), and that of \(B'\) is equal to B (VIII, p.~259, Theorem 5). The relation \(L \supset B'\) therefore implies \(L \subset B\) . We have proved a).

Let us prove b). Suppose that B contains \(B'\) , and choose a maximal semisimple commutative subalgebra L of A contained in B, which exists by a). Let g be a homomorphism from B to A. By Proposition 3 of VIII, p.~262, we have the equality \([A:K]=[L:K]^{2}\) ; by the corollary of VIII, p.~263, there exists an inner automorphism \(\theta_{1}\) of A that coincides with g on L. If f is the canonical injection of B into A, then the homomorphisms g and \(\theta_{1}\circ f\) have the same restriction to the center \(B'\) of B because \(B'\) is contained in L. By Theorem 2 of VIII, p.~256, there exists an inner automorphism \(\theta\) of A such that \(g=\theta\circ f\) ; in other words, \(\theta\) extends g.

